\documentclass[10pt,a4paper]{amsart}
\usepackage[margin=19mm]{geometry}
\usepackage{amsmath,amssymb,mathtools}
\usepackage[scr=rsfs,cal=euler]{mathalfa}
\usepackage{xfrac}
\usepackage[unicode,hidelinks,bookmarks=false]{hyperref}
\newcommand{\ie}{i.e.\ }
\newcommand{\cf}{cf.\ }
\newcommand{\resp}{respectively\ }
\newcommand{\Subsection}{\subsection}
\providecommand{\nobreakdash}{\nobreak\hbox{-}}
\setlength{\emergencystretch}{2em}
\multlinegap0pt
\usepackage{amssymb}
\usepackage{mathtools}
\newtheorem{lemma}{Lemma}
\newcommand{\RR}{\mathbb{R}}
\title{Persistence of periodic billiard orbits under domain deformation}
\author{Samuel Everett}
\date{Source: arXiv:2605.04362v2 (CC BY 4.0)}
\begin{document}
\maketitle

\noindent\emph{Source and licence.} Persistence of periodic billiard orbits under domain deformation. Version arXiv:2605.04362v2 (CC BY 4.0), distributed by arXiv under the Creative Commons Attribution 4.0 International licence, http://creativecommons.org/licenses/by/4.0/, as recorded in the arXiv licence metadata for this exact version and shown on the abstract page https://arxiv.org/abs/2605.04362v2. Full licence notice: \texttt{licenses/arxiv-2605.04362-CC-BY-4.0.txt}.\par\smallskip

\noindent\textit{Editorial scope.} Counted unit: the article's lemma lem:contFP (source0.tex lines 284--286). The article's complete original proof of it is delivered with it (source0.tex lines 287--304, one \texttt{proof} environment of 18 lines). No mathematics is written, repaired or re-derived: the statement and the proof are the article's own lines, in the article's own notation, and no retained line is substituted or re-flowed.  No further source-local context is needed: every cross-reference of every retained line resolves inside the two delivered spans.  Nothing was omitted from any retained span, so the package carries no omission record.  No retained line contains a citation, so the wrapper carries no bibliography and no bibliography entry.  No retained line contains a figure environment, an image-inclusion command or drawing code, so no image is delivered and the package declares no asset dependency.  Editorial formatting only, in addition: an AMS \texttt{amsart} preamble that restates the article's own declarations and macros used by the retained lines, namely the environments \texttt{lemma} on separate counters, together with the standard packages those lines need.  The retained environments print in this document as Lemma 1 in order of appearance, whereas the article prints the same items with the numbering of its own full document.  The complete native source of the article as distributed in the arXiv e-print for this exact version is bundled unchanged as \texttt{sources/199853/source0.tex}.
\begin{lemma}\label{lem:contFP}
Let $P = (0,A)^n\subset \RR^n$. Suppose $f:P\times \RR\rightarrow \RR$ is a jointly continuous function such that for each $p \in P$, $x\mapsto f(p, x)$ is a contraction on $\RR$. Then the unique fixed point $x^*(p) \in \RR$ is continuous with respect to $p$ on $P$.
\end{lemma}

\begin{proof}
Let $p \in P$ be arbitrary. The function $f_p: \mathbb{R} \to \mathbb{R}$ defined by $f_p(x) = f(p, x)$ is a jointly continuous contraction mapping on the complete metric space $\mathbb{R}$, with a contraction constant $L=L(p) < 1$. By the contraction mapping theorem, there is a unique $x^*(p) \in \mathbb{R}$ such that $f(p, x^*(p)) = x^*(p)$. So, the mapping $x^*: P \to \mathbb{R}$ is well-defined.
Define $G:P\times\mathbb R\to\mathbb R$ by $G(p,x):=f(p,x)-x$. Then $G$ is also jointly continuous.

Take $p\in(0,1)^n$. If $x>y$, and $L(p)\in[0,1)$ is a contraction constant for $f_p$, then
\begin{align*}
G(p,x)-G(p,y)
&=\bigl(f(p,x)-f(p,y)\bigr)-(x-y)\\
&\le |f(p,x)-f(p,y)|-(x-y)\\
&\le (L(p)-1)(x-y)\\
&<0,
\end{align*}
so the function $x\mapsto G(p,x)$ is strictly decreasing and continuous. In particular, since $G(p,x^*(p))=0$ and the zero is unique, it follows that $G(p,x)>0$ for $x<x^*(p)$, and $G(p,x)<0$ for $x>x^*(p)$.

To prove continuity of $x^*$, fix $p_0\in P$ and set $x_0=x^*(p_0)$. Given $\varepsilon>0$, let $a = x_0-\varepsilon$ and $b=x_0+\varepsilon$. By the foregoing sign characterization at $p_0$, we have $G(p_0,a)>0$ and $G(p_0,b)<0$. By joint continuity of $G$, there exists $\delta>0$ such that $\|p-p_0\|<\delta$ implies $G(p,a)>0$ and $G(p,b)<0$.

For such $p$, continuity of $x\mapsto G(p,x)$ and the intermediate value theorem yields the existence of some $x\in(a,b)$ with $G(p,x)=0$. Uniqueness of the zero (equivalently, uniqueness of the fixed point of the contraction $f_p$) forces $x=x^*(p)$, hence $|x^*(p)-x_0|<\varepsilon$. Since $\varepsilon>0$ was arbitrary, $x^*$ is continuous at $p_0$, and because $p_0$ was arbitrary, the map $p\mapsto x^*(p)$ is continuous on $P$.
\end{proof}

\end{document}
